% ---- Font settings ------
\usepackage{luatexja-fontspec}
\setmainfont{Noto Serif}
\setmainjfont{Noto Serif CJK JP}

% --- Margin ---
\usepackage[margin=25truemm]{geometry}

% --- TOC ---
\usepackage[hidelinks]{hyperref}
\usepackage{tocloft}
\renewcommand{\cftsecleader}{\cftdotfill{\cftdotsep}}
\setlength{\cftsecnumwidth}{4em}

% --- Between the line ---
\usepackage{setspace}
\setstretch{1.15}

% --- Heading --- 
\usepackage{titlesec}
\titlespacing{\section}{0pt}{2em}{0.8em}

% --- Formula ---
\usepackage{amsmath, amssymb}

% --- Figure ---
\usepackage{graphicx}

% --- Code ----
\usepackage{xcolor}
\usepackage{float}
\usepackage{caption}
\usepackage{newfloat}
\usepackage{etoolbox}
\usepackage{minted}
\setminted{
  bgcolor=black!3,
  frame=lines,
  framesep=1mm,
  fontsize=\footnotesize,
  linenos=false,
  breaklines=true,
  autogobble=true
}
\renewcommand{\listingscaption}{ソースコード}

% 端末出力用のフロート環境（「出力 1」と表示される）
\DeclareFloatingEnvironment[fileext=lop, listname=出力一覧]{termout}
\SetupFloatingEnvironment{termout}{name=出力}

% キャプション設定: 上置き
\captionsetup[termout]{position=top, skip=0pt, aboveskip=0pt, belowskip=6pt, labelsep=colon}
\captionsetup[listing]{position=top, skip=0pt, aboveskip=0pt, belowskip=0pt, labelsep=colon}

% [H] フロートの前後の空き（既定 12pt）を詰める
\setlength{\intextsep}{2pt plus 2pt minus 2pt}

% 出力ブロック内の minted が足す余分な縦空きを消す
\AtBeginEnvironment{termout}{\setlength{\topsep}{0pt}\setlength{\partopsep}{0pt}}
\AtBeginEnvironment{listing}{\setlength{\topsep}{0pt}\setlength{\partopsep}{0pt}}

% --- Bibliography ---
\usepackage{tocbibind}

% --- URL ---
\usepackage{url}

% --- Assignment Env ---
\newcommand{\courcetitle}{情報理工学実験II}
\newcommand{\reporttitle}{ネットワークプログラミング}
\newcommand{\affiliation}{工学部 情報エレクトロニクス学科}
\newcommand{\studentyear}{学士3年}
\newcommand{\studentnum}{02246910}
\newcommand{\faculitynum}{ICT26903}
\newcommand{\studentname}{吉野 拓人}
\newcommand{\mailaddress}{takuto.yoshino.l3@elms.hokudai.ac.jp}

\usepackage{fancyhdr}
\pagestyle{fancy}
\fancyhf{}
\rhead{\reporttitle}
\cfoot{\thepage}

% --- env for separating by assignment ---
\newenvironment{maintask}[1]{
  \section*{課題 #1}
  \addcontentsline{toc}{section}{課題 #1}
}{}

\newenvironment{extask}[1]{
  \section*{発展課題 #1}
  \addcontentsline{toc}{section}{発展課題 #1}
}{}

% --- Custom Title ---
\newcommand{\makecustomtitle}{
  \begin{titlepage}
    \begin{center}
      \vspace*{30mm} % 上部の余白
      
      % 講義名とレポート回数
      {\Large \courcetitle{}} \\[2em]
      {\huge \textbf{\reporttitle}} \\
      
      \vspace{60mm} % 中央の大きな空白が「感覚」を生みます
      
      % 個人情報セクション（左寄せに近い中央配置）
      \begin{minipage}{0.8\textwidth}
        \centering
        \large
        \begin{tabular}{clp{10cm}}
          学部学科  &: \quad \affiliation \\[0.8em]
          学年      &: \quad \studentyear \\[0.8em]
          学生番号  &: \quad \studentnum \\[0.8em]
          学部番号  &: \quad \faculitynum \\[0.8em]
          学生氏名  &: \quad \studentname \\[0.8em]
          メールアドレス &: \quad \mailaddress \\[0.8em]
        \end{tabular}
      \end{minipage}

      \vfill % 残りのスペースをすべて埋める
      
      % 提出日
      {\large 提出日 ： \number\year 年 \number\month 月 \number\day 日}
      
      \vspace*{20mm}
    \end{center}
  \end{titlepage}
  \clearpage
}
